\subsection{Síntesis de integración con el Caso 10}\label{suxedntesis-de-integraciuxf3n-con-el-caso-10}

La conexión entre el marco regulatorio y la operación real de Transportes Curimón S.A. demostró que las exigencias legales no admiten una revisión superficial. Al cruzar el marco normativo (P2 y P3), el presupuesto (P4) y la arquitectura tecnológica (P5), aparecen tres relaciones importantes:

Primero, el régimen de corresponsabilidad del Art. 15 bis de la Ley N.º 21.719 obliga a distinguir operativamente entre los 196 conductores propios (amparados en el contrato de trabajo y el Dictamen DT 569/2018) y los 258 conductores dependientes de 148 transportistas subcontratados. Para estos últimos, la licitud del rastreo telemático exige acuerdos de encargo (DPA) específicos y cláusulas de consentimiento expreso en la aplicación móvil antes de liberar turnos.

Segundo, la falta de una Agencia de Protección de Datos constituida y de Cláusulas Contractuales Tipo oficiales en Chile genera una zona gris respecto a la réplica de contingencia en Azure East US 2 y el tránsito a Mendoza. AudIT aborda esta brecha usando de forma preventiva las Cláusulas Tipo del RGPD europeo (Capítulo V) y dejando consideradas en el TCO horas de asesoría legal para ajustarlas cuando la autoridad emita directrices.

Tercero, cada control técnico exigido por las bases queda conectado con una partida concreta del presupuesto: el cifrado de base de datos a nivel de campo (RT-11.10) se cubre con Azure Key Vault (725,4 UF), la gestión del RAT se apoya en la plataforma GRC CISO Assistant Pro (255 UF) y el cumplimiento de ISO/IEC 27001 se respalda con consultoría E-26 y auditoría externa acreditada (387,5 UF). Así se evita proponer controles sin costo asociado o gastos sin justificación normativa.

\subsection{Conclusiones y recomendaciones estratégicas}\label{conclusiones-y-recomendaciones-estratuxe9gicas}

\subsection{Conclusiones}\label{conclusiones}

\begin{enumerate}
\def\labelenumi{\arabic{enumi}.}
\tightlist
\item
  \textbf{El cumplimiento es necesario para operar:} La investigación dejó en claro que cumplir con las leyes no es solo un trámite legal o papeleo, sino algo indispensable para que el proyecto funcione en la práctica. Si no se cuenta con bases legales claras o no se avisa a tiempo de incidentes, Curimón se arriesga a multas muy graves que pueden llegar a 20.000 UTM por datos personales o 40.000 UTM por ciberseguridad.
\item
  \textbf{Solución a la desorganización de datos:} El problema de fondo en Curimón era que la información de los camiones y choferes estaba repartida en planillas Excel sin ningún orden ni seguridad. La propuesta de conectar la flota por capas y unificarla en una sola vista resuelve este desorden de raíz, asegurando que los datos viajen protegidos y que no se pierda información en las zonas sin señal.
\item
  \textbf{Inversión justificada:} Gastar \textbf{8.492,16 UF netas} (\(\text{VAN}_{\text{costo}} = \mathbf{6.651,24\text{ UF}}\)) representa apenas el 3,9\% de lo que cuesta toda la licitación. Con solo un 2,80\% de probabilidad al año de recibir una fiscalización con sanción, el plan de seguridad ya se paga completamente solo, logrando un retorno sobre la inversión (RoSI) superior al \(+250\%\).
\end{enumerate}

\subsection{Recomendaciones estratégicas priorizadas}\label{recomendaciones-estratuxe9gicas-priorizadas}

\begin{enumerate}
\def\labelenumi{\arabic{enumi}.}
\tightlist
\item
  \textbf{Fase 1 · Ordenar contratos y permisos (Meses 1 a 4):} Lo primero es firmar los contratos de encargo con los 148 transportistas externos y pedir el consentimiento en la app a los 258 choferes subcontratados. Además, se debe armar el registro de datos (RAT) antes de activar el rastreo masivo por GPS.
\item
  \textbf{Fase 2 · Montar la plataforma y el enlace de seguridad (Meses 5 a 12):} Desplegar los sistemas en Azure Chile Central, configurar las 686 claves en Key Vault para cifrar los datos de cada transportista y dejar listo el canal de alertas para avisarle al CSIRT Nacional en menos de 3 horas si ocurre un ataque.
\item
  \textbf{Fase 3 · Certificación formal y respaldo final (Meses 13 a 20):} Hacer la auditoría externa para certificar la norma ISO/IEC 27001:2022 con una empresa acreditada y contratar el seguro de ciberriesgo con Chubb, dejando el proyecto completamente respaldado frente a los 84 clientes corporativos.
\end{enumerate}
